\section{Безразмерная постановка задачи}
	
	Для приведения системы уравнений Эйлера к безразмерному виду используем $\pi$-теорему теории размерности. В качестве основных масштабов выберем:
	\begin{itemize}
		\item Масштаб длины: $L = R$ — радиус цилиндра,
		\item Масштаб скорости: $U_\infty$ — скорость набегающего потока,
		\item Масштаб плотности: $\rho_\infty$ — плотность невозмущенного потока,
		\item Масштаб давления: $\rho_\infty U_\infty^2$.
	\end{itemize}
	
	Введем следующие безразмерные переменные:
	\[
	x = \frac{x^*}{L}, \quad y = \frac{y^*}{L}, \quad \varphi = \frac{\varphi^*}{U_\infty L}, \quad \rho = \frac{\rho^*}{\rho_\infty}, \quad p = \frac{p^*}{\rho_\infty U_\infty^2} \\
	\quad u = \frac{u^*}{U_\infty}, \quad a = \frac{a^*}{a_\infty}\quad M_{loc} = \frac{q}{a}\cdot M_\infty.
	\]
	
	Размер расчетной области задан величиной $R_\infty = 20$. Основным варьируемым параметром задачи является число Маха набегающего потока $M_\infty$.
	
	\subsection{Безразмерные уравнения Эйлера}
	
	\subsubsection{Уравнение неразрывности}
	\[
	\frac{\partial (\rho u)}{\partial x} + \frac{\partial (\rho v)}{\partial y} = 0.
	\]
	
	\subsubsection{Уравнения движения (компоненты)}
	\[
	\rho \left( u \frac{\partial u}{\partial x} + v \frac{\partial u}{\partial y} \right) = -\frac{\partial p}{\partial x},
	\]
	\[
	\rho \left( u \frac{\partial v}{\partial x} + v \frac{\partial v}{\partial y} \right) = -\frac{\partial p}{\partial y}.
	\]
	
	\subsubsection{Уравнение для потенциала скорости}
	При условии безвихревости течения $\nabla \times \vec{v} = 0$ существует потенциал скорости $\vec{v} = \nabla \varphi$. Уравнение неразрывности в безразмерной форме принимает вид:
	\[
	\frac{\partial}{\partial x} \left( \rho \frac{\partial \varphi}{\partial x} \right) + \frac{\partial}{\partial y} \left( \rho \frac{\partial \varphi}{\partial y} \right) = 0.
	\]
	
	\subsubsection{Уравнение изэнтропичности и связь плотности с потенциалом}
	
	Из интеграла Бернулли для безвихревого изэнтропического течения:
	\[
	\frac{\gamma}{\gamma-1} \frac{p}{\rho} + \frac{1}{2} (u^2 + v^2) = \frac{1}{(\gamma-1)M^2_\infty} + \frac{1}{2}.
	\]
	
	С учетом изэнтропического соотношения $p / \rho^\gamma = \text{const}$ и переходя к безразмерным переменным, получим:
	\[
	\rho = \left[ 1 + \frac{\gamma-1}{2} M_\infty^2 \left( 1 - (u^2 + v^2) \right) \right]^{\frac{1}{\gamma-1}},
	\]
	где $M_\infty = U_\infty / a_\infty$ — число Маха набегающего потока, $a_\infty = \sqrt{\gamma p_\infty / \rho_\infty}$ — скорость звука в невозмущенном потоке.
	
	В терминах потенциала скорости:
	\[
	u = \frac{\partial \varphi}{\partial x}, \quad v = \frac{\partial \varphi}{\partial y},
	\]
	\[
	q^2 = u^2 + v^2 = \left( \frac{\partial \varphi}{\partial x} \right)^2 + \left( \frac{\partial \varphi}{\partial y} \right)^2.
	\]
	
	Таким образом, плотность выражается через потенциал скорости:
	\begin{equation}
		\boxed{\rho = \left[ 1 + \frac{\gamma-1}{2} M_\infty^2 \left( 1 - \left( \frac{\partial \varphi}{\partial x} \right)^2 - \left( \frac{\partial \varphi}{\partial y} \right)^2 \right) \right]^{\frac{1}{\gamma-1}}}
		\label{eq:density}
	\end{equation}
	
	\subsection{Безразмерные граничные условия}
	
	\subsubsection{Условие непротекания на поверхности цилиндра}
	На поверхности цилиндра $r = 1$ (в безразмерных координатах) нормальная компонента скорости равна нулю:
	\[
	\vec{v} \cdot \vec{n} = 0 \quad \text{при} \quad r = 1.
	\]
	В полярных координатах это соответствует условию:
	\[
	\frac{\partial \varphi}{\partial r} = 0 \quad \text{при} \quad r = 1.
	\]
	
	\subsubsection{Условие на бесконечности}
	На больших расстояниях от цилиндра поток полагается невозмущенным (условие Дирихле):
	\[
	\varphi \to x = r \cos \theta \quad \text{при} \quad r \to \infty.
	\]
	
	
	
	%%%%%%%%%%%%%%%%%%%%%%%%%%%%%%%%%
